% Source: https://artofproblemsolving.com/community/c3h1213086_ratio_of_areas
% Stable identity: aops-topic-1213086-q01
The diagonal of the square is a diameter, of length $2r$, so the square has side $r\sqrt2$. Each triangle is an isosceles right triangle with legs $r\sqrt2$ and hypotenuse $2r$.

\begin{asy}
usepackage("amsmath");
size(3.5cm);
defaultpen(fontsize(10pt));
real r=1; pair A=(r,0),B=(0,r),C=(-r,0),D=(0,-r);
draw(circle((0,0),r)); draw(A--B--C--D--cycle,linewidth(0.9)); draw(A--C,blue);
real q=r*(sqrt(2)-1); pair I=(0,q); draw(circle(I,q),red); dot(I);
label("$r$",(0.5,0),S); label("$r\sqrt2$",(A+B)/2,NE);
label("\mbox{}",(min(currentpicture,user=true)+max(currentpicture,user=true))/2,black); // keeps the text after the figure black when LaTeX typesets the labels inline
\end{asy}

The inradius of a right triangle is $\dfrac{\text{leg}+\text{leg}-\text{hypotenuse}}{2}$, so
\[\rho=\dfrac{2r\sqrt2-2r}{2}=r(\sqrt2-1).\]
The ratio of areas is $\left(\dfrac{\rho}{r}\right)^2=(\sqrt2-1)^2=\boxed{3-2\sqrt2}\approx0.172$.
